
\subsubsection{Research Questions}

Three hypotheses were tested:

\begin{itemize}
\item \textbf{H-E1 (Existence):} Can a linear classifier predict training benchmark from penultimate layer representations with accuracy > 60\%?
\item \textbf{H-M1 (Mechanism):} Does BFS correlate with cross-dataset gap (r > 0.3, p < 0.05)?
\item \textbf{H-M2 (Mechanism):} Do single-benchmark models show larger gaps than multi-benchmark models (> 5 percentage points difference)?
\end{itemize}

\subsubsection{Datasets}

\begin{table*}[htbp]
\centering
\resizebox{\dimexpr 2\columnwidth+\columnsep\relax}{!}{%
\begin{tabular}{lllll}
\toprule
Dataset & Classes & Domain & Train Size & Test Size \\
\midrule
Flowers102 & 102 & Fine-grained flowers & 2,040 & 6,149 \\
CIFAR-100 & 100 & Coarse-grained objects & 50,000 & 10,000 \\
\bottomrule
\end{tabular}
}
\end{table*}

The original experimental design specified five fine-grained benchmarks (CUB-200, Stanford Dogs, Stanford Cars, FGVC Aircraft, Oxford Flowers) with NABirds as a held-out evaluation set. Due to scope constraints, the proof-of-concept execution used two benchmarks (Flowers102, CIFAR-100), which increases chance level from 20\% to 50\%.

\subsubsection{Model Configuration}

\begin{itemize}
\item 6 models total: 2 benchmarks $\times$ 3 random seeds
\item Architecture: ResNet-50 pretrained on ImageNet-1K
\item Fine-tuning: 10 epochs, SGD, lr=0.01, momentum=0.9, cosine annealing
\item Feature dimension: 2048 (avgpool layer)
\end{itemize}

\subsubsection{Probe Configuration}

\begin{itemize}
\item Classifier: Logistic regression (sklearn default parameters)
\item Training data: 20,000 feature vectors (5,000 per model $\times$ 4 models, using CIFAR-100 test images)
\item Test data: 10,000 feature vectors (5,000 per model $\times$ 2 held-out models)
\end{itemize}

